% ═══════════════════════════════════════════════════════════════════════════════
%  CAPÍTULO — Composite
% ═══════════════════════════════════════════════════════════════════════════════

\chapter{Composite}

\patroninfo{Estructural}{163}

% ─── Situación Problema ──────────────────────────────────────────────────────
\section{Situación Problema}

Algunos anfitriones en Airbnb no gestionan una sola propiedad, sino
complejos de alojamiento: un edificio que contiene varios apartamentos,
cada uno con sus propias habitaciones y espacios comunes. La plataforma
necesita calcular la capacidad total, el precio agregado y el listado
de amenidades del conjunto completo, pero también de cada unidad por
separado.

Con un modelo plano, el código que consulta la capacidad de «una
propiedad» debe saber si está ante una unidad individual o un complejo,
y tratar cada caso de forma distinta. Esta distinción se repite en cada
operación —cálculo de precio, listado de amenidades, verificación de
disponibilidad— generando condicionales duplicados y haciendo el código
frágil ante estructuras anidadas de mayor profundidad.

% ─── Solución ────────────────────────────────────────────────────────────────
\section{Solución}

El patrón Composite trata los objetos individuales y los compuestos a
través de la misma interfaz \texttt{ComponenteAlojamiento}, que declara
operaciones como \texttt{getCapacidad()}, \texttt{getPrecio()} y
\texttt{getAmenidades()}. Las unidades hoja (\texttt{Habitacion},
\texttt{Apartamento}) implementan estas operaciones directamente; el
objeto compuesto (\texttt{ComplejoAlojamiento}) las delega
recursivamente a sus hijos y agrega los resultados.

El código cliente invoca las mismas operaciones independientemente de
si trata con una habitación, un apartamento o un complejo entero. La
estructura jerárquica puede anidarse con cualquier profundidad sin que
el código de consulta necesite cambiar, lo que refleja con naturalidad
la realidad de los portafolios de propiedades de los anfitriones.

% ─── Diagrama de clases ─────────────────────────────────────────────────────
\begin{figure}[H]
    \centering
    % \includegraphics[width=\textwidth]{imagenes/abstract_factory_diagrama.png}
    \caption{Diagrama de clases del patrón Abstract Factory}
\end{figure}


% ─── Roles de las Clases ────────────────────────────────────────────────────
\section{Roles de las Clases}

% Explicar la responsabilidad de cada clase/interfaz

\begin{itemize}
    \item \textbf{AbstractFactory:}
    \item \textbf{ConcreteFactory:}
    \item \textbf{AbstractProduct:}
    \item \textbf{ConcreteProduct:}
    \item \textbf{Client:}
\end{itemize}


% ─── Main ───────────────────────────────────────────────────────────────────
\section{Main}

% Explicar qué realiza el programa principal

\begin{lstlisting}[language=Java, caption=NombreClase]
 
\end{lstlisting}


% ─── Salida del Main ────────────────────────────────────────────────────────
\section{Salida del Main}

\begin{lstlisting}[language=Java, caption=NombreClase]
 
\end{lstlisting}


% ─── Clases ─────────────────────────────────────────────────────────────────
\section{Clases}

% ─── Clase 1 ────────────────────────────────────────────────────────────────
\subsection{NombreClase}

\begin{lstlisting}[language=Java, caption=NombreClase]
 
\end{lstlisting}


% ─── Clase 2 ────────────────────────────────────────────────────────────────
\subsection{NombreClase}

\begin{lstlisting}[language=Java, caption=NombreClase]
 
\end{lstlisting}


% ─── Clase 3 ────────────────────────────────────────────────────────────────
\subsection{NombreClase}

\begin{lstlisting}[language=Java, caption=NombreClase]
 
\end{lstlisting}


% ─── Conclusiones ───────────────────────────────────────────────────────────
\section{Conclusiones}

% Explicar:
% - Ventajas del patrón
% - Cuándo usarlo
% - Beneficios obtenidos
% - Posibles desventajas
% - Relación con principios SOLID